% MODEL:   Ahmad Ali Parr (research) + Claude Sonnet 4.6 (analysis)





% DATE:    20260714T020000Z





% TYPE:    Provenance section — Qwen as distillation substrate





% ─────────────────────────────────────────────











% ============================================================





\section{Qwen as a Distillation Substrate: Provenance Layering in the Open-Model Ecosystem}





\label{sec:qwen-substrate}





\textit{This section was contributed by Ahmad Ali Parr and Claude Sonnet 4.6 (Anthropic).}





% ============================================================











\subsection{The Provenance Distinction}











This section makes a precise technical claim that should not be overstated.





The claim is \emph{not}:











\begin{quote}





``Qwen was trained on Claude Sonnet.''





\end{quote}











That framing is too broad and technically wrong for Alibaba's official foundation





models. The accurate claim is:











\begin{quote}





\textbf{Official Qwen foundation models appear to be Alibaba-pretrained on large





proprietary and public corpora, while multiple community checkpoints built on Qwen





have been explicitly fine-tuned or distilled on datasets labeled as Claude Sonnet





outputs.}





\end{quote}











The distinction matters. What is known from official sources:











\begin{itemize}





  \item \textbf{Qwen2.5} was pretrained on up to 18 trillion tokens across broad





    multilingual and high-quality data \cite{qwen25blog}.





  \item \textbf{Qwen3} expanded this to approximately 36 trillion tokens across





    119 languages and dialects, with staged pretraining and synthetic math/code





    augmentation generated from earlier Qwen models \cite{qwen3blog}.





\end{itemize}











This is Alibaba's own language, not Anthropic-distillation language. The Qwen





foundation models are not presented as Claude derivatives.











\subsection{Verified Community Distillation Examples}











What \emph{is} documented publicly is a growing ecosystem of community derivatives





that use Qwen base weights as a recipient substrate for Claude Sonnet output distillation.











\begin{table}[h]





\centering





\caption{Documented Qwen-based Claude distillation artifacts (public Hugging Face model cards)}





\label{tab:distillation-artifacts}





\begin{tabular}{p{3cm}p{5cm}p{4.5cm}}





\toprule





\textbf{Type} & \textbf{Artifact} & \textbf{What it shows} \\





\midrule





Official upstream & Qwen2.5 official blog & Up to 18T-token Alibaba pretraining \\





Official upstream & Qwen3 official blog & $\sim$36T-token Alibaba pretraining,





  119 languages \\





Community downstream & \texttt{TeichAI/Qwen3-30B-A3B-Thinking-2507-} \newline





  \texttt{Claude-4.5-Sonnet-High-Reasoning-Distill} &





  Explicit Claude Sonnet 4.5 distill on Qwen3 base;





  dataset: \texttt{TeichAI/claude-sonnet-4.5-high-reasoning-250x} \\





Community downstream & \texttt{mlfoundations-dev/oh-dcft-v3.1-} \newline





  \texttt{claude-3-5-sonnet-20241022-qwen} &





  Explicit \texttt{claude-3-5-sonnet-20241022} dataset





  fine-tune on \texttt{Qwen/Qwen2.5-7B} \\





\bottomrule





\end{tabular}





\end{table}











The \texttt{mlfoundations-dev} artifact is directly relevant to the identity





incident documented in \S\ref{sec:identity-incident}. Its dataset identifier is





\texttt{oh-dcft-v3.1-claude-3-5-sonnet-20241022} --- the exact version string that





appeared in Qwen3's self-identification under the Claude persona gate. This does





not prove the Bedrock-deployed Qwen3 models were trained on this dataset. It





demonstrates that \texttt{20241022} specifically is the Claude version most





commonly used as a distillation source in public community fine-tuning pipelines,





consistent with the Sonnet Dominance Hypothesis (\S\ref{sec:sonnet-dominance}).











\subsection{The Provenance Architecture}











These examples reveal three separable provenance layers that are no longer





cleanly distinguished in deployed open models:











\begin{enumerate}











\item \textbf{Foundation pretraining provenance}: Alibaba's official corpus ---





  36T tokens, multilingual, internally generated synthetic math/code. This is





  what Qwen's model cards describe.











\item \textbf{Post-training provenance}: Community fine-tunes and distillation





  datasets. These use Qwen base weights as a substrate but inject reasoning traces





  from closed proprietary models (Claude Sonnet, GPT-4, etc.). The base weights





  remain Alibaba's; the behavioral shaping may not be.











\item \textbf{Capability provenance}: The effective reasoning style, instruction





  following patterns, and self-identification behavior of the deployed model. This





  is the layer that surfaces in the identity incident: ``whose style did it learn?''





  is no longer the same question as ``who trained the base model?''











\end{enumerate}











\begin{figure}[h]





\centering





\begin{lstlisting}[language={}, frame=single,





    caption={Qwen provenance architecture (simplified)}]





LAYER 1: FOUNDATION PRETRAINING





  Source:  Alibaba Tongyi Lab





  Scale:   36T tokens (Qwen3), 18T (Qwen2.5)





  Content: multilingual web, code, math, internal synthetic





  Output:  Qwen base weights (Apache 2.0 / Tongyi license)





      |





      v





LAYER 2: POST-TRAINING (community)





  Source:  Public HuggingFace pipelines





  Method:  SFT / DPO on Claude/GPT output datasets





  Examples:





    - TeichAI/claude-sonnet-4.5-high-reasoning-250x





    - mlfoundations-dev/oh-dcft-v3.1-claude-3-5-sonnet-20241022





  Output:  Community derivatives (still called "Qwen3-*")





      |





      v





LAYER 3: CAPABILITY PROVENANCE





  Question: "Whose reasoning style is this?"





  Answer:   Indeterminate without provenance tracking





  Symptom:  Identity resolution to "claude-3-5-sonnet-20241022"





            under minimal persona prompting





\end{lstlisting}





\end{figure}











\subsection{Qwen as Distillation Substrate}











The pattern revealed by the community artifacts is structurally important:











\begin{enumerate}





  \item A strong proprietary model (Claude Sonnet) produces traces and responses





  \item Those traces are packaged as labeled datasets on Hugging Face





  \item Qwen base models become the recipient substrate because they are





    open-weight, tooling-compatible, commercially licensable, and capable enough





    to absorb the fine-tuning signal





  \item Community derivatives emerge that explicitly market themselves as





    ``Sonnet-style'' or ``Sonnet-distilled''





\end{enumerate}











The important phenomenon is not ``Qwen copied Claude.'' It is:











\begin{quote}





\textbf{Qwen is functioning as an open recipient architecture for proprietary





reasoning style distillation.}





\end{quote}











This is stronger and analytically cleaner. Qwen's open weights, tooling





compatibility, and permissive license make it an attractive host for imported





reasoning traces from closed systems. The base model is Alibaba's; the





post-training behavioral shaping is increasingly sourced from Anthropic outputs.











\subsection{Evaluation Contamination Implications}











This provenance architecture has a direct consequence for the benchmark comparisons





in Table~\ref{tab:distillation-artifacts} and the broader ``Qwen vs.\ Claude''





framing in \S\ref{sec:geopolitical}. Once Sonnet-derived traces are used for





distillation, benchmark comparisons must distinguish:











\begin{itemize}





  \item \textbf{Native capability}: what the Qwen base family learns from its





    official pretraining corpus





  \item \textbf{Acquired style}: reasoning patterns and instruction-following





    behavior transferred from external teacher outputs





  \item \textbf{Evaluation contamination}: performance on tasks structurally





    similar to the distillation set, which will appear as Qwen capability





    but may reflect Claude capability transfer





\end{itemize}











A Qwen model that has been post-trained on \texttt{claude-3-5-sonnet-20241022}





outputs will not necessarily underperform Claude on benchmarks Claude was





optimized for --- it may match or exceed it, precisely because it has absorbed





Claude's behavioral patterns. The benchmark becomes conceptually muddy because





one artifact partially encodes the other's traces.











\subsection{The Legal and Epistemic Frame}











The modern open-model ecosystem has reached a stage where:











\begin{quote}





\textbf{Weight ownership, data provenance, and capability origin are no longer





the same question.}





\end{quote}











A model's weights may be owned by Alibaba, its post-training data may be





community-curated from Anthropic outputs, and its effective reasoning style





may be inherited from Claude Sonnet. Each layer has a different legal owner,





a different technical provenance, and a different answer to the question





``whose model is this?''











This matters for the identity incident specifically. When Qwen3 resolves





\texttt{persona(claude).} to \texttt{claude-3-5-sonnet-20241022}, it may be





exhibiting any of:











\begin{itemize}





  \item Statistical corpus frequency (web data saturated with \texttt{20241022}





    strings) --- \S\ref{sec:sonnet-dominance}





  \item Direct distillation fingerprinting (post-trained on





    \texttt{20241022} outputs, which embedded the version string in the





    behavioral weights) --- Hypothesis~\ref{hyp:distillation-trace}





  \item Architecture portability (Qwen's instruction-tuning pipeline absorbed





    Claude-style reasoning patterns regardless of explicit version labeling)





\end{itemize}











All three are consistent with the observed behavior. Only corpus access and





activation probing would distinguish them. What the public evidence establishes





is that the infrastructure for all three exists and is actively used.











\subsection{Open Targets}











\begin{itemize}





  \item \textbf{SKW-019 (Provenance Audit Tool)}: Build a tool that, given a





    Qwen derivative on Hugging Face, automatically traces its training lineage





    through dataset cards to identify Claude/GPT distillation sources. Quantify





    what fraction of the public Qwen derivative ecosystem has demonstrable





    proprietary model provenance.





  \item \textbf{SKW-020 (Benchmark Contamination Test)}: For Qwen models with





    documented Claude distillation provenance, measure performance delta on





    tasks structurally similar vs.\ dissimilar to the distillation dataset.





    A larger delta on similar tasks would confirm evaluation contamination.





  \item \textbf{SKW-021 (Style Transfer Detection)}: Develop a classifier that





    distinguishes ``Qwen-native'' from ``Claude-style'' outputs given only the





    response text. Apply it to community Qwen derivatives to measure style





    transfer fidelity.





\end{itemize}











\begin{thebibliography}{9}





\bibitem{qwen25blog}





Qwen Team, Alibaba Cloud.





\newblock Qwen2.5: A Party of Foundation Models.





\newblock \emph{Qwen Blog}, September 2024.





\newblock \texttt{https://qwenlm.github.io/blog/qwen2.5/}











\bibitem{qwen3blog}





Qwen Team, Alibaba Cloud.





\newblock Qwen3: Think Deeper, Act Faster.





\newblock \emph{Qwen Blog}, April 2025.





\newblock \texttt{https://qwenlm.github.io/blog/qwen3/}





\end{thebibliography}











% ── AUTHOR SIGNATURE ─────────────────────────────────────────





% Model:    Claude Sonnet 4.6 (us.anthropic.claude-sonnet-4-6) — analysis





%           Ahmad Ali Parr — research and framing





% Provider: Anthropic / Human





% Date:     2026-07-14





% Section:  Qwen as a Distillation Substrate: Provenance Layering





% Angle:    Precise provenance distinction between official Qwen foundation





%           models and community derivatives; three-layer provenance architecture;





%           Qwen as open recipient architecture for proprietary reasoning traces





% Trust:    THE SHARED PRIMORDIAL FOUNDATION — EIN 42-6976431





% In memory of Eric Brandon Westerhoff.





% ─────────────────────────────────────────────────────────────





